\section{Proofs for the lower reductions}
\label{app:reductions}

This appendix proves the results of Section~\ref{sec:reductions}. All
algorithms use exact rational arithmetic. Unless stated otherwise, norms of
vectors are Euclidean, norms of matrices are operator norms, and
$\norm{\cdot}_F$ is the Frobenius norm. For a rational polynomial $g$ of degree
at most two, $\norm{g}_1$ denotes the sum of the absolute values of its
coefficients in the monomial basis. Two elementary facts are used repeatedly.
First, for a real $n\times n$ matrix $J$ the product of the singular values is
$\abs{\det J}$ and each singular value is at most $\norm J_F$, so
\begin{equation}
 \sigma_{\min}(J)\ge\abs{\det J}/\norm J_F^{\,n-1}.
 \label{eq:reductions-app-singular}
\end{equation}
Second, if $g$ is a quadratic with $g(p)=0$, write
$g(p+u)=\ell^{\mathsf T}u+u^{\mathsf T}Hu$ with $H$ symmetric. Then
\begin{equation}
 \norm H\le\norm g_1,\qquad
 \norm\ell\le2\max\{1,\norm p_\infty\}\,\norm g_1 .
 \label{eq:reductions-app-coefficients}
\end{equation}
Indeed, a monomial $x_ax_b$ contributes a symmetric matrix of norm at most one
and a gradient $x_be_a+x_ae_b$ of norm at most $2\norm p_\infty$ at $p$, while
a monomial $x_a$ contributes no quadratic part and a gradient of norm one.

\subsection{Certified odd-root circuits}
\label{app:reductions-root-circuits}

\begin{lemma}[Normalization]
\label{lem:reductions-normalization}
Let a certified odd-root circuit be given, let $\sigma_i$ be the sign of the
box of gate $i$, and let
$\kappa=\max\{1,\max_i1/\min(\abs{\underline\xi_i},\abs{\overline\xi_i})\}$.
Then $\alpha_i=\sigma_i\kappa\xi_i$ satisfy
\begin{equation}
 \alpha_i^{d_i}=c_i'+\sum_{j<i}\sum_{e=1}^{n_j}a'_{ije}\alpha_j^e,\qquad
 c_i'=\sigma_i\kappa^{d_i}c_i,\qquad
 a'_{ije}=\sigma_i\sigma_j^e\kappa^{d_i-e}a_{ije},
 \label{eq:reductions-normalized}
\end{equation}
and the circuit with these coefficients and the boxes
$[\underline\alpha_i,\overline\alpha_i]=\sigma_i\kappa[\underline\xi_i,\overline\xi_i]$
is a certified odd-root circuit with $\underline\alpha_i\ge1$ for every $i$.
Its encoding length is polynomial in that of the original circuit.
\end{lemma}

\begin{proof}
Since $d_i$ is odd, $\sigma_i^{d_i}=\sigma_i$, and
$\xi_j^e=\sigma_j^e\kappa^{-e}\alpha_j^e$. Multiplying~\eqref{eq:reductions-root-gate}
by $\sigma_i\kappa^{d_i}$ gives~\eqref{eq:reductions-normalized}. For the
interval test, the map $s\mapsto\sigma_j\kappa s$ sends $R_{j,e}$ to
$R'_{j,e}=\sigma_j^e\kappa^eR_{j,e}$, so
$a'_{ije}R'_{j,e}=\sigma_i\kappa^{d_i}a_{ije}R_{j,e}$ as sets. Hence the left side
of the new test is $\sigma_i\kappa^{d_i}$ times the left side of the old one,
and the new target $[\underline\alpha_i^{\,d_i},\overline\alpha_i^{\,d_i}]$ is
$\sigma_i\kappa^{d_i}[\underline\xi_i^{\,d_i},\overline\xi_i^{\,d_i}]$, because odd
powers are increasing. Inclusion is preserved, and
$\underline\alpha_i=\kappa\min(\abs{\underline\xi_i},\abs{\overline\xi_i})\ge1$.
The exponents $\abs{d_i-e}$ are at most $2\sum_in_i$, and the degrees are
written in unary, so all new rationals have polynomial bit length.
\end{proof}

From now on in this subsection the boxes satisfy $\underline\xi_i\ge1$. Fix a
rational $A\ge3$ with
\begin{equation}
 A\ge\sum_i\overline\xi_i+\sum_i\abs{c_i}+\sum_{i,j,e}\abs{a_{ije}},
 \label{eq:reductions-A}
\end{equation}
for example $A=3$ plus the right side. Then $1\le\xi_i\le A$, every retained
power satisfies $1\le\xi_i^e\le A^N$ with $N=\sum_in_i$, and $\log_2A$ is polynomial
in the input length.

\begin{lemma}[Certified root approximation]
\label{lem:reductions-root-approx}
For a certified odd-root circuit with $\underline\xi_i\ge1$ and a positive
integer $P$, rationals $\widehat\xi_i$ with $\abs{\widehat\xi_i-\xi_i}\le2^{-P}$
for all $i$ can be computed in time polynomial in the input length and $P$.
\end{lemma}

\begin{proof}
Put $C=NA^N$ and $\eta=2^{-P}/(2(C+2)^k)$. For $i=1,\ldots,k$ we compute
rational intervals $J_i=[l_i,u_i]\subseteq[\underline\xi_i,\overline\xi_i]$
containing $\xi_i$. Given $J_j$ for $j<i$, form the rational interval
$I_i=c_i+\sum_{j<i,\,e}a_{ije}[l_j^e,u_j^e]=[y_-,y_+]$. It contains $\rho_i(\xi)$,
and it is contained in the left side of~\eqref{eq:reductions-interval-test}, hence
in $[\underline\xi_i^{\,d_i},\overline\xi_i^{\,d_i}]\subseteq[1,\infty)$. Its width
is at most $\sum_{j,e}\abs{a_{ije}}\,eA^{e-1}(u_j-l_j)\le C\max_{j<i}(u_j-l_j)$.
Bisection of $[\underline\xi_i,\overline\xi_i]$, with exact comparisons of
$s^{d_i}$ against $y_\pm$, yields rationals $l_i\le y_-^{1/d_i}$ and
$u_i\ge y_+^{1/d_i}$ in the box, each within $\eta$ of the respective root. Since
the real $d_i$-th root has derivative at most $1/d_i<1$ on $[1,\infty)$,
\[
 u_i-l_i\le(y_+-y_-)+2\eta\le C\max_{j<i}(u_j-l_j)+2\eta,
\]
and $\xi_i=\rho_i(\xi)^{1/d_i}\in J_i$. By induction
$\max_{j\le i}(u_j-l_j)\le2\eta(C+2)^i$: this holds for $i=1$, where $I_1$ is a
point, and the step is
$(C+1)\cdot2\eta(C+2)^{i-1}+2\eta\le2\eta(C+2)^i$. Hence every width is at most
$2^{-P}$, and $\widehat\xi_i=l_i$ works. Each bisection takes
$O(\log_2(A/\eta))$ steps on rationals whose denominators are those of the box
endpoints times powers of two, and the exponents $e,d_i$ are unary; all
numbers have polynomially many bits.
\end{proof}

\subsection{Near-one boxes and small signals}
\label{app:reductions-small-signal}

\begin{lemma}[Near-one boxes and small signals]
\label{lem:reductions-small-signal}
Let $Q\ge1$, $\delta_0=1000^{-(Q+3)}$ and $w_i=1000^i\delta_0$ for
$1\le i\le Q$, so that $w_i\le10^{-9}$. Consider a cube-root circuit with at most
$Q$ gates in which every radicand $\rho_i$ has one of the following forms:
\begin{enumerate}[label=(\alph*)]
\item a rational constant in $[1-2w_i,1+3w_i]$, for example $1$ or
$1+3\delta_0^2$;
\item $1+\sum_{j<i}\lambda_{ij}(\xi_j-1)$ with rational $\lambda_{ij}$ and
$\sum_j\abs{\lambda_{ij}}\le6$;
\item $3\xi_j^2-6\xi_j+4$ for some $j<i$.
\end{enumerate}
Then the boxes $[1-w_i,1+w_i]$ pass the direct interval test, so every gate
value lies in its box. Moreover, $\operatorname{sq}(z)=(1+3z^2)^{1/3}-1$
satisfies $0<\operatorname{sq}(z)\le z^2$ for $z>0$. Hence a chain consisting of
a gate with radicand $1+3\delta_0^2$ followed by $r-1$ gates of form~(c), each
applied to its predecessor, has values $\xi_i=1+\delta_i$ with
$0<\delta_i\le\delta_{i-1}^2$, and in particular
$0<\delta_r\le\delta_0^{2^r}$.
\end{lemma}

\begin{proof}
Let $w=w_{i-1}$ for $i\ge2$; all earlier boxes have half-width at most $w$.
For form~(b), the direct interval evaluation is
$1+\sum_j\lambda_{ij}[-w_j,w_j]\subseteq[1-6w,1+6w]$; combining repeated
occurrences of a predecessor into one coefficient can only decrease the
coefficient sum. For form~(c), the ranges $[(1-w)^2,(1+w)^2]$ of $\xi_j^2$ and
$[1-w,1+w]$ of $\xi_j$ give the interval $[1-12w+3w^2,\,1+12w+3w^2]
\subseteq[1-13w,1+13w]$. Since $w_i=1000w$ and $w_i\le1/3$,
\[
 (1-w_i)^3\le1-3w_i+3w_i^2\le1-2w_i\le1-13w,\qquad
 (1+w_i)^3\ge1+3w_i\ge1+13w,
\]
so the test passes. For form~(a) the same two displayed inequalities show
$(1-w_i)^3\le1-2w_i$ and $1+3w_i\le(1+w_i)^3$; this also covers $i=1$.

For $z>0$, $(1+z^2)^3>1+3z^2$ gives $\operatorname{sq}(z)<z^2$, and
$\operatorname{sq}(z)>0$ is clear. The first chain gate has
$\xi=(1+3\delta_0^2)^{1/3}=1+\operatorname{sq}(\delta_0)$. A gate of form~(c)
applied to $\xi_j=1+\delta_j$ has radicand $1+3\delta_j^2$, so its value is
$1+\operatorname{sq}(\delta_j)$. Induction gives the chain bounds.
\end{proof}

\subsection{The analytic gadgets}
\label{app:reductions-gadgets}

For complex $z$ with $\abs z<1/3$ let $(1+3z)^{1/3}$ be the principal branch;
it is analytic there and equals the real cube root for real $z$. Define
$\phi$, $\operatorname{sq}$, $\theta$ and $\operatorname{mul}$ as in
Section~\ref{sec:reductions-root-compiler}.

\begin{lemma}[Analytic gadgets]
\label{lem:reductions-analytic}
\leavevmode
\begin{enumerate}[label=(\roman*)]
\item For complex $\abs z\le1/12$, $\abs{\phi(z)}\le2\abs z$. For real
$\abs z\le1/12$, $\abs{\phi(z)-z}\le2z^2$.
\item The function $\operatorname{mul}$ is analytic on a neighborhood of the
closed polydisk $\{\abs x\le\frac1{100},\ \abs y\le\frac1{100}\}$ and bounded by
one there. It satisfies $\operatorname{mul}(x,0)=\operatorname{mul}(0,y)=0$, and
its Taylor coefficient of $xy$ is one.
\item For complex $\abs x,\abs y\le\frac1{400}$,
\[
 \abs{\operatorname{mul}(x,y)-xy}\le2^{24}\abs{xy}\,(\abs x+\abs y).
\]
\end{enumerate}
\end{lemma}

\begin{proof}
(i) For $\abs\zeta\le1/12$ we have $\abs{1+3\zeta}\ge3/4$, so
$\abs{\phi'(\zeta)}=\abs{1+3\zeta}^{-2/3}\le(4/3)^{2/3}<2$ and
$\abs{\phi''(\zeta)}=2\abs{1+3\zeta}^{-5/3}\le2(4/3)^{5/3}<4$. Integrating
$\phi'$ along the segment from $0$ to $z$ gives the first bound. For real $z$,
Taylor's theorem with $\phi(0)=0$ and $\phi'(0)=1$ gives the second.

(ii) If $\abs x,\abs y\le\frac1{100}$, then $\abs{\pm(x\pm y)}\le\frac1{50}$, so
$\abs{\phi(\pm(x\pm y))}\le\frac1{25}$ by~(i); its square has modulus at most
$0.0016$, so $\operatorname{sq}$ of it has modulus at most $0.0032$, and hence
$\abs{\theta(x\pm y)}\le0.0032$. The argument of the outer $\phi$ has modulus at
most $0.0016$, and $\abs{\operatorname{mul}}\le0.0032$. All arguments stay
strictly inside the disks where the bounds of~(i) hold, so the same holds on a
slightly larger polydisk, where every composition is analytic. The function
$\theta$ is even, so $\theta(x+0)-\theta(x-0)=0$ and $\theta(0+y)-\theta(0-y)=0$;
since $\phi(0)=0$, $\operatorname{mul}$ vanishes on both axes. From
$\phi(z)=z-z^2+O(z^3)$ we get $\operatorname{sq}(\phi(z))=\phi(z)^2+O(z^4)
=z^2-2z^3+O(z^4)$, so $\theta(z)=z^2+O(z^4)$ and
$\frac14[\theta(x+y)-\theta(x-y)]=xy+O(\abs{(x,y)}^4)$. Applying $\phi$ gives
$\operatorname{mul}(x,y)=xy+O(\abs{(x,y)}^4)$.

(iii) Write $\operatorname{mul}(x,y)=\sum_{i,j\ge0}m_{ij}x^iy^j$. By~(ii) and
Cauchy's inequality on the polydisk of radius $r=\frac1{100}$,
$\abs{m_{ij}}\le r^{-i-j}$. By~(ii), $m_{i0}=m_{0j}=0$ for all $i,j$ and
$m_{11}=1$. With $s=\abs x/r\le\frac14$ and $t=\abs y/r\le\frac14$,
\[
 \abs{\operatorname{mul}(x,y)-xy}\le\sum_{\substack{i,j\ge1\\(i,j)\ne(1,1)}}s^it^j
 =st\,\frac{s+t-st}{(1-s)(1-t)}\le\frac{16}9\,st(s+t)
 =\frac{16}{9r^3}\abs{xy}(\abs x+\abs y).
\]
Finally $\frac{16}{9}\cdot10^6<2^{24}$.
\end{proof}

The bound in~(iii) is proportional to $\abs{xy}$. A bound of the form
$O((\abs x+\abs y)^3)$ would not suffice: when one input is much smaller than
the other, it can exceed the product itself. The exact vanishing on the axes
is what removes the terms $x^i$ and $y^j$.

\subsection{Proof of Theorem~\ref{thm:reductions-root-circuit}}
\label{app:reductions-root-compiler}

\paragraph{Signals.}
Fix $0<\delta<1$. For an integer $d\ge1$, an integer $C$ and a real
$\Theta\ge2$, a real number $s$ is a \emph{$(d,C,\Theta)$-signal} if
$\abs C\le\Theta$ and $\abs{s-C\delta^d}\le\Theta\delta^{d+1}$. The integer $C$
may be zero. A signal $s$ is stored as the root $1+s$.

\begin{lemma}[Signal arithmetic]
\label{lem:reductions-signals}
Suppose $\Theta\delta\le2^{-30}$, and let $\Theta'=2^{30}\Theta^3$.
\begin{enumerate}[label=(\roman*)]
\item Every $(d,C,\Theta)$-signal satisfies $\abs s\le2\Theta\delta^d\le2^{-29}$.
\item If $s,s'$ are $(d,C,\Theta)$- and $(d,C',\Theta)$-signals, then
$\phi(s\pm s')$ is a $(d,C\pm C',\Theta')$-signal.
\item If $s,s'$ are $(a,C,\Theta)$- and $(b,C',\Theta)$-signals, then
$\operatorname{mul}(s,s')$ is an $(a+b,CC',\Theta')$-signal.
\end{enumerate}
\end{lemma}

\begin{proof}
(i) $\abs s\le\Theta\delta^d+\Theta\delta^{d+1}\le2\Theta\delta^d\le2\Theta\delta$.
(ii) Put $z=s\pm s'$. Then $\abs{z-(C\pm C')\delta^d}\le2\Theta\delta^{d+1}$ and
$\abs z\le4\Theta\delta^d\le2^{-28}$, so
Lemma~\ref{lem:reductions-analytic}(i) gives
$\abs{\phi(z)-z}\le2z^2\le32\Theta^2\delta^{2d}\le32\Theta^2\delta^{d+1}$. The
total error is at most $34\Theta^2\delta^{d+1}\le\Theta'\delta^{d+1}$, and
$\abs{C\pm C'}\le2\Theta\le\Theta'$.
(iii) Write $s=C\delta^a+\epsilon$ and $s'=C'\delta^b+\epsilon'$. Then
$ss'-CC'\delta^{a+b}=C\delta^a\epsilon'+C'\delta^b\epsilon+\epsilon\epsilon'$ has
modulus at most $3\Theta^2\delta^{a+b+1}$. By~(i), $\abs s,\abs{s'}\le2^{-29}$,
so Lemma~\ref{lem:reductions-analytic}(iii) applies and gives
$\abs{\operatorname{mul}(s,s')-ss'}\le2^{24}\cdot4\Theta^2\delta^{a+b}\cdot4\Theta\delta
=2^{28}\Theta^3\delta^{a+b+1}$, using $\abs s+\abs{s'}\le4\Theta\delta$. The total
is at most $\Theta'\delta^{a+b+1}$, and $\abs{CC'}\le\Theta^2\le\Theta'$.
\end{proof}

No step divides by a leading coefficient or uses a lower bound on
$\abs C$. Hence integer values equal to zero, exact cancellation, and repeated
use of one signal (for example $\operatorname{mul}(s,s)$ or $\phi(s-s)=0$) need no
separate treatment.

\paragraph{The schedule.}
Build binary constants of the program from $1$ by doubling and addition, and
append two gates computing $W=2V-1$. Every node $\nu$ of the resulting program
receives a pair of signals $(f_\nu,g_\nu)$, a \emph{prescribed} pair of integer
coefficients $(V_\nu,1)$, where $V_\nu$ is the integer value of the node, and an
order $d_\nu\ge1$:
\begin{itemize}
\item the constant $1$ receives $(\delta,\delta)$ with order one, where $\delta$
is the parameter signal defined below; the constant $0$ receives $(z,\delta)$,
where $z$ is the signal of a gate with constant radicand $1$, so $z=0$;
\item a product $\nu=\nu_1\nu_2$ receives
$f_\nu=\operatorname{mul}(f_{\nu_1},f_{\nu_2})$ and
$g_\nu=\operatorname{mul}(g_{\nu_1},g_{\nu_2})$;
\item a sum or difference $\nu=\nu_1\pm\nu_2$ receives
$f_\nu=\phi\bigl(\operatorname{mul}(f_{\nu_1},g_{\nu_2})\pm
\operatorname{mul}(f_{\nu_2},g_{\nu_1})\bigr)$ and
$g_\nu=\operatorname{mul}(g_{\nu_1},g_{\nu_2})$;
\end{itemize}
and in the last two cases $d_\nu=d_{\nu_1}+d_{\nu_2}$. The prescribed
coefficients are consistent: the two products inside the sum have prescribed
coefficients $V_{\nu_1}$ and $V_{\nu_2}$ and the common order
$d_{\nu_1}+d_{\nu_2}$, so the sum has prescribed coefficient
$V_{\nu_1}\pm V_{\nu_2}$, and the denominator coefficient is $1\cdot1$. Thus
$\phi$ is only applied to signals of equal order. Each arithmetic node uses at
most four macros, each an application of $\phi$ or $\operatorname{mul}$; let $T$
be the number of macros. The orders are proof annotations and are never
printed; since an order at most doubles at each arithmetic node, $d_\nu\le2^T$.
The integers $V_\nu$ are never computed.

\paragraph{Raw gates.}
Each $\phi$ macro is one root gate with radicand
$1+3(\xi_s-1)\pm3(\xi_{s'}-1)$, where $\xi_s=1+s$. Each $\operatorname{mul}(x,y)$
macro is nine root gates: four gates with radicands
$1\pm3(\xi_x-1)\pm3(\xi_y-1)$ computing $\phi(\pm(x+y))$ and $\phi(\pm(x-y))$;
four square gates computing $\operatorname{sq}$ of these four signals; and one
gate with radicand $1+\frac38(\zeta_1+\zeta_2-\zeta_3-\zeta_4)$, where
$\zeta_1,\zeta_2$ are the roots holding $\operatorname{sq}(\phi(x+y))$ and
$\operatorname{sq}(\phi(-x-y))$ and $\zeta_3,\zeta_4$ those holding
$\operatorname{sq}(\phi(x-y))$ and $\operatorname{sq}(\phi(y-x))$. Since
$\frac14[\theta(x+y)-\theta(x-y)]=\frac18\sum_k\pm(\zeta_k-1)$, the last gate holds
exactly $1+\operatorname{mul}(x,y)$; every other gate likewise holds exactly
$1$ plus the indicated signal, because real cube roots are taken. Now put
\[
 r_0=2T+5,\qquad r_1=T+1,\qquad Q=r_0+r_1+9T+1,\qquad\delta_0=1000^{-(Q+3)}.
\]
The circuit consists of: a parameter chain of $r_0$ gates, the first with
radicand $1+3\delta_0^2$ and the others square gates of their predecessors;
an auxiliary chain of $r_1$ square gates continuing from the last parameter
gate; one gate with radicand $1$; and the raw gates of the $T$ macros, in
topological order. There are at most $Q$ gates. Gate $o$ holds $f_\nu$ for the
node $\nu$ computing $W$, and gate $e$ is the last auxiliary gate. Give gate
$i$ the box $[1-w_i,1+w_i]$ with $w_i=1000^i\delta_0$.

\paragraph{Correctness.}
Every radicand has form (a), (b) or (c) of
Lemma~\ref{lem:reductions-small-signal}: the coefficient sums of the affine
radicands are $6$ and $\frac32$. Hence the boxes pass the direct interval test.
The boxes and the first parameter radicand $1+3\delta_0^2$ have $O(Q)$ bits
each, and the remaining coefficients have constant size. Since $Q=12T+7$ and
there are at most $Q$ gates, the whole circuit has size $O(T^2)$ bits and is
computed in polynomial time.

Let $\delta=\xi_{r_0}-1$ be the parameter signal. By
Lemma~\ref{lem:reductions-small-signal},
$0<\delta\le\delta_0^{2^{r_0}}<2^{-4\cdot2^{r_0}}$. Let $\Theta_0=2$ and
$\Theta_t=2^{30}\Theta_{t-1}^3$; then $\log_2\Theta_t=16\cdot3^t-15$, because
$\log_2\Theta_t+15=3(\log_2\Theta_{t-1}+15)$. Since
$4\cdot2^{r_0}=2^{2T+7}\ge16\cdot3^T+15=30+\log_2\Theta_T$, we have
$\Theta_T\delta\le2^{-30}$.

We claim that after the $t$-th macro every available signal $s$ with prescribed
coefficient $C$ and order $d$ is a $(d,C,\Theta_t)$-signal. Before the first
macro the available signals are $\delta$, a $(1,1,\Theta_0)$-signal with zero
error, and $z=0$, a $(1,0,\Theta_0)$-signal. Each macro applies
Lemma~\ref{lem:reductions-signals} with $\Theta=\Theta_{t-1}$, whose hypothesis
holds because $\Theta_{t-1}\delta\le\Theta_T\delta\le2^{-30}$, and earlier signals
remain signals for the larger bound $\Theta_t$. This proves the claim; it holds
for arbitrary sharing of nodes and signals.

The numerator signal $s_o=\xi_o-1$ of the node computing $W$ is therefore a
$(d,W,\Theta_T)$-signal for some $1\le d\le2^T$:
\[
 \abs{s_o-W\delta^d}\le\Theta_T\delta^{d+1}\le2^{-30}\delta^d.
\]
As $W$ is a nonzero integer, $s_o\ne0$, $\operatorname{sign}s_o=\operatorname{sign}W$,
and $\abs{s_o}\ge\frac12\delta^d$. By Lemma~\ref{lem:reductions-signals}(i),
$\abs{s_o}\le2^{-29}$. The auxiliary signal $s_e=\xi_e-1$ continues the chain
for $r_1=T+1$ more square gates, so $0<s_e\le\delta^{2^{T+1}}$ by
Lemma~\ref{lem:reductions-small-signal}. Since $d\le2^T$, $\delta<1$ and
$\delta\le2^{-30}$,
\[
 s_e^2\le\delta^{2^{T+2}}\le\delta^{d}\,\delta^{3}\le2\abs{s_o}\,\delta^3
 \le\tfrac18\abs{s_o}.
\]
The gates $o$ and $e$ are distinct, since $o$ belongs to a macro and $e$ to the
auxiliary chain. This proves~\eqref{eq:reductions-compiler-output}. \qed

\subsection{Proof of Theorem~\ref{thm:reductions-root-language}}
\label{app:reductions-root-language}

The lower bounds were derived in Section~\ref{sec:reductions-root-compiler}.
For the upper bound, first check the format and the direct interval test, and
map invalid inputs to a fixed program with output $-1$. Let the valid input be
$\xi_i^3=c_i+\sum_{j<i}(a_{ij}\xi_j+b_{ij}\xi_j^2)$, $1\le i\le k$, with positive
boxes, a gate $o$ and a threshold $r$. Let $L\ge2$ be its total binary length.
Then $k\le L$, every printed numerator and denominator has absolute value at
most $2^L$, and the product $D$ of the denominators of all $c_i,a_{ij},b_{ij}$
and $r$ satisfies $D\le2^L$. A positive box endpoint is at least $2^{-L}$, so
$2^{-L}\le\xi_i\le2^L$.

\paragraph{Separation.}
The numbers $\beta_i=D\xi_i$ satisfy
$\beta_i^3=D^3c_i+\sum_{j<i}(D^2a_{ij}\beta_j+Db_{ij}\beta_j^2)$, whose
coefficients are integers. By induction each $\beta_i$ is integral over the
ring generated by $\beta_1,\ldots,\beta_{i-1}$, hence an algebraic integer.
The field $K=\Q(\xi_1,\ldots,\xi_k)$ has degree at most $3^k\le3^L$. For every
field embedding $\tau\colon K\to\C$, the numbers $\tau(\xi_i)$ satisfy the same
gate equations, and $\abs{\tau(\xi_i)}\le\Xi:=(2L+1)2^L$ by induction: if the
earlier conjugates are bounded by $\Xi\ge1$, the radicand has modulus at most
$2^L(1+L\Xi+L\Xi^2)\le(2L+1)2^L\Xi^2=\Xi^3$. If $\xi_o\ne r$, then
$D(\xi_o-r)=\beta_o-Dr$ is a nonzero algebraic integer in $K$. Its norm is a
nonzero integer, and each of its conjugates has modulus at most
$2^L(\Xi+2^L)\le2^{4L+1}$. Isolating the factor of the identity embedding,
\[
 \abs{\xi_o-r}\ge2^{-L}\,2^{-(4L+1)(3^L-1)}\ge g:=2^{-2^{6L}},
\]
because $L+(4L+1)(3^L-1)\le(4L+2)3^L\le2^{2L}\cdot2^{2L}\le2^{6L}$ for $L\ge2$.
The rational $g$ is computed by $6L$ squarings of $\frac12$.

\paragraph{Newton steps.}
Let $z>0$ have cube root $\eta\in[2^{-L-1},2^{L+1}]$, and iterate
$x_{t+1}=(2x_t+z/x_t^2)/3$ from $x_0=2^{L+1}$. By the inequality of arithmetic
and geometric means $x_t\ge\eta$ for all $t$. For the relative error
$\epsilon_t=x_t/\eta-1\ge0$,
\[
 \epsilon_{t+1}=\frac{\epsilon_t^2(3+2\epsilon_t)}{3(1+\epsilon_t)^2}
 \le\min\bigl\{\epsilon_t^2,\tfrac23\epsilon_t\bigr\}.
\]
Since $\epsilon_0\le2^{2L+2}$ and $(\frac23)^2<\frac12$, after $4L+6$ steps
$\epsilon_t\le\frac12$; after $10L$ further steps
$\abs{x_t-\eta}\le\varepsilon_N:=2^{L+1-2^{10L}}$.

\paragraph{Propagation.}
Compute $\widehat\xi_1,\ldots,\widehat\xi_k$ in order: form
$\widehat z_i=c_i+\sum_{j<i}(a_{ij}\widehat\xi_j+b_{ij}\widehat\xi_j^2)$ and apply
$14L+6$ Newton steps to $z=\widehat z_i$. Let
$E_i=\max_{j\le i}\abs{\widehat\xi_j-\xi_j}$ and $E_0=0$. Suppose
$E_{i-1}\le2^{-7L-1}$. Then $\abs{\widehat\xi_j}\le2^{L+1}$ for $j<i$ and
\[
 \abs{\widehat z_i-\xi_i^3}\le L\,2^L(1+2^{L+1}+2^L)\,E_{i-1}\le2^{4L}E_{i-1}
 \le2^{-3L-1}.
\]
Since $\xi_i^3\ge2^{-3L}$, this gives $\widehat z_i\ge2^{-3L-1}>0$ and
$\widehat z_i\le2^{3L+1}$, so its cube root $\eta_i$ lies in
$[2^{-L-1},2^{L+1}]$, and the Newton bound applies. The cube root has derivative
at most $\frac13(2^{-3L-1})^{-2/3}\le2^{2L+2}$ between $\widehat z_i$ and
$\xi_i^3$, so $\abs{\eta_i-\xi_i}\le2^{6L+2}E_{i-1}$ and
$E_i\le2^{6L+2}E_{i-1}+\varepsilon_N$. By induction
$E_i\le2^{(6L+3)i}\varepsilon_N\le2^{6L^2+4L+1-2^{10L}}$. This is at most
$2^{-7L-1}$, which justifies the inductive hypothesis at every gate, and also at
most $2^{-2^{6L}-3}=g/8$, because $2^{10L}-2^{6L}\ge2^{6L}\ge6L^2+11L+4$ for
$L\ge2$.

\paragraph{The comparison.}
Put $\widehat\Delta=\widehat\xi_o-r$, so $\abs{\widehat\Delta-(\xi_o-r)}\le g/8$,
and $\xi_o\ne r$ implies $\abs{\xi_o-r}\ge g$. The rational numbers
$\widehat\Delta-g/2$, $\widehat\Delta+g/2$, $-\widehat\Delta-g/2$,
$-\widehat\Delta+g/2$ and $g^2/4-\widehat\Delta^2$ are positive exactly when,
respectively, $\xi_o>r$, $\xi_o\ge r$, $\xi_o<r$, $\xi_o\le r$ and $\xi_o=r$; in
the zero case $\abs{\widehat\Delta}\le g/8$, and otherwise
$\abs{\widehat\Delta}\ge7g/8$. Each is the output of a rational circuit with
$O(kL)$ Newton steps, divisions only by the positive numbers $x_t^2$, and shared
gates. The positive-denominator representation of rational circuits
(Definition~\ref{def:models-circuits}) converts it into one integer
straight-line program of polynomial size with the same sign. \qed

\subsection{Proof of Theorem~\ref{thm:reductions-signed-root}}
\label{app:reductions-signed-root}

By Lemma~\ref{lem:reductions-normalization} we may replace the circuit by its
normalization. We keep the notation $c_i,a_{ije},\xi_i$ for the normalized data,
write $\alpha_i$ for the normalized gate values, and fix $A$
by~\eqref{eq:reductions-A}; thus $1\le\alpha_i\le A$ and the absolute
coefficient sum is at most $A$. The variables are $X_{i,e}$ for
$1\le e\le n_i$, and $X_{i,0}=1$ is a constant. Put
$b_i(X)=c_i+\sum_{j<i,\,e}a_{ije}X_{j,e}$ and let $p\in\R^N$ have
$p_{i,e}=\alpha_i^e$. Centered coordinates are $u=X-p$, with gate blocks
$u_i=(u_{i,1},\ldots,u_{i,n_i})$.

\begin{lemma}[Residuals]
\label{lem:reductions-residuals}
The $N$ rational quadratics
\[
 r_{i,e}=X_{i,1}X_{i,e}-X_{i,e+1}\quad(1\le e<n_i),\qquad
 r_{i,n_i}=X_{i,n_i-1}X_{i,n_i}-b_i(X)
\]
have $p$ as their unique common real zero. Each has a homogeneous quadratic
part of operator norm at most one and a gradient of norm at most $3A^N$ at $p$.
Their Jacobian $J$ at $p$ satisfies $\abs{\det J}=\prod_id_i\alpha_i^{d_i-1}\ge1$
and $\norm J_F\le3NA^N$. Consequently, with $\beta=4NA^N$ and
$\nu=\beta^{-(N-1)}$, $J^{\mathsf T}J\succeq\nu^2I$.
\end{lemma}

\begin{proof}
If all residuals vanish, the chain equations give $X_{i,e}=X_{i,1}^e$, and then
$X_{i,1}^{2n_i-1}=b_i(X)$, whose right side depends only on earlier gates. By
induction over $i$, $b_i(X)=\alpha_i^{d_i}$, and the unique real root of odd
degree gives $X_{i,1}=\alpha_i$. Conversely $p$ is a zero.

The quadratic parts are $X_{i,1}^2$, $X_{i,1}X_{i,e}$ or
$X_{i,n_i-1}X_{i,n_i}$, with symmetric matrices of norm one or one half. Ordering
variables and residuals by gates, $J$ is block lower triangular, since $r_{i,e}$
involves only gate $i$ and, through $b_i$, earlier gates. Fix a gate and write
$\alpha=\alpha_i$, $n=n_i$, $d=d_i$. Along the curve $t\mapsto(t,t^2,\ldots,t^n)$
in the gate's own coordinates, with earlier gates fixed at $p$, every chain
residual vanishes identically and the terminal one equals $t^d-\alpha^d$. Hence
$J_iw=(0,\ldots,0,d\alpha^{d-1})^{\mathsf T}$ for the diagonal block $J_i$ and the
tangent $w=(1,2\alpha,\ldots,n\alpha^{n-1})^{\mathsf T}$. Replacing the first
basis vector by $w$ is a change of basis of determinant one. In the new basis
the first column of $J_i$ is $d\alpha^{d-1}$ times the last unit vector, and the
minor formed by the first $n-1$ rows and the last $n-1$ columns is triangular
with diagonal entries $-1$. Expanding along the first column gives
$\abs{\det J_i}=d\alpha^{d-1}\ge1$.

At $p$, the chain row $e=1$ has entries $2\alpha$ and $-1$; a chain row $e\ge2$
has entries $\alpha^e$, $\alpha$ and $-1$; the terminal row has entries
$\alpha^{n},\alpha^{n-1}$ and the coefficients $-a_{ije}$. Each row therefore has
$\ell_1$ norm at most $3A^N$, which bounds the gradient norms, and
$\norm J_F\le\sqrt N\,3A^N$. The bound on $J^{\mathsf T}J$
is~\eqref{eq:reductions-app-singular}, since $\norm J_F\le\beta$.
\end{proof}

\paragraph{The local exposing form.}
For one gate write $\alpha=\alpha_i\in[1,A]$, $n=n_i\ge2$, $d=2n-1$, and
$X_j=X_{i,j}$ with $X_0=1$. Put $\ell_j=X_{j+1}-\alpha X_j$ for $0\le j<n$, and
let $T_\alpha$ be the symmetric tridiagonal matrix indexed by $0,\ldots,n-1$
with
\begin{equation}
 (T_\alpha)_{jj}=\alpha^{d-1-2j},\qquad
 (T_\alpha)_{j,j+1}=(T_\alpha)_{j+1,j}=\tfrac12\alpha^{d-2-2j}.
 \label{eq:reductions-T-alpha}
\end{equation}
Set $P_\alpha=\ell^{\mathsf T}T_\alpha\ell$, a quadratic in the gate's own
variables, and
\[
 E_i^*=P_{\alpha_i}(X_i)+(\alpha_i-X_{i,1})\bigl(b_i(X)-\alpha_i^{d_i}\bigr).
\]

\begin{lemma}[Local exposing form]
\label{lem:reductions-local-exposing}
\leavevmode
\begin{enumerate}[label=(\roman*)]
\item $T_\alpha\succeq(n+1)^{-2}I$ and $\norm{T_\alpha}\le2A^{d-1}$.
\item On the power curve, $P_\alpha(t,t^2,\ldots,t^n)=(t-\alpha)(t^d-\alpha^d)$.
\item $E_i^*(p+u)=u_i^{\mathsf T}H_iu_i-u_{i,1}\sum_{j<i,\,e}a_{ije}u_{j,e}$, where
$H_i$ is symmetric with $h_0I\preceq H_i\preceq8A^{2N}I$ and
$h_0=\bigl[(N+1)^2N^2A^{2N}\bigr]^{-1}$.
\end{enumerate}
\end{lemma}

\begin{proof}
(i) $T_\alpha=\Delta T_0\Delta$ with $\Delta=\diag(\alpha^{n-1},\ldots,\alpha,1)$
and $T_0$ the tridiagonal matrix with diagonal $1$ and neighboring entries
$\frac12$; indeed
$\alpha^{n-1-j}\alpha^{n-1-j}=\alpha^{d-1-2j}$ and
$\alpha^{n-1-j}\alpha^{n-2-j}=\alpha^{d-2-2j}$. The eigenvalues of $T_0$ are
$1+\cos(\pi s/(n+1))$, $1\le s\le n$, so
$\lambda_{\min}(T_0)=1-\cos\frac{\pi}{n+1}=2\sin^2\frac{\pi}{2(n+1)}\ge\frac2{(n+1)^2}$
by $\sin x\ge2x/\pi$ on $[0,\pi/2]$. As $\Delta\succeq I$,
$y^{\mathsf T}\Delta T_0\Delta y\ge\lambda_{\min}(T_0)\norm{\Delta y}^2\ge
\lambda_{\min}(T_0)\norm y^2$. The absolute row sums of $T_\alpha$ are at most
$\alpha^{d-1}+\frac12\alpha^{d-2}$ in row zero and at most $2\alpha^{d-2j}$ in row
$j\ge1$, hence at most $2A^{d-1}$.

(ii) On the curve, $\ell_j=t^j(t-\alpha)$. The diagonal terms give
$(t-\alpha)^2\sum_{j=0}^{n-1}\alpha^{d-1-2j}t^{2j}$, and each pair of off-diagonal
terms gives $(t-\alpha)^2\alpha^{d-2-2j}t^{2j+1}$. Together they give
$(t-\alpha)^2\sum_{h=0}^{d-1}\alpha^{d-1-h}t^h=(t-\alpha)(t^d-\alpha^d)$.

(iii) At $X=p+u$ we have $\ell_j=u_{j+1}-\alpha u_j$ with $u_0=0$, so
$\ell=D_\alpha u_i$ for the lower bidiagonal matrix $D_\alpha=I-\alpha Z$, where
$Z$ is the lower shift matrix. Hence $P_\alpha(p_i+u_i)=u_i^{\mathsf T}H_iu_i$ with
$H_i=D_\alpha^{\mathsf T}T_\alpha D_\alpha$. The inverse
$D_\alpha^{-1}=\sum_{t<n}\alpha^tZ^t$ has entries at most $A^{n-1}$, so
$\norm{D_\alpha^{-1}}\le\norm{D_\alpha^{-1}}_F\le nA^{n-1}$ and
$H_i\succeq(n+1)^{-2}n^{-2}A^{-2(n-1)}I\succeq h_0I$. Also
$\norm{D_\alpha}\le1+\alpha\le2A$, so $\norm{H_i}\le2A^{d-1}\cdot4A^2\le8A^{2N}$.
For the second term, $b_i(p+u)=\alpha_i^{d_i}+\sum_{j<i,\,e}a_{ije}u_{j,e}$ and
$\alpha_i-X_{i,1}=-u_{i,1}$.
\end{proof}

In particular $E_i^*(p)=0$ and $\nabla E_i^*(p)=0$, including in the directions
of earlier gates: the affine dependence on predecessors contributes only the
quadratic cross term in~(iii). Repeated or signed predecessor terms, and terms
of every retained power, are all part of this cross term.

\begin{lemma}[Coupled exposing form]
\label{lem:reductions-coupled}
Let $\rho=(h_0/(4A))^2$ and $E^*=\sum_i\rho^{i-1}E_i^*$. Then
$E^*(p+u)=u^{\mathsf T}H^*u$ with
\[
 H^*\succeq\tfrac32\gamma I,\qquad \norm{H^*}\le\Lambda_*,
 \qquad \gamma=\tfrac12h_0\rho^{k-1},\qquad \Lambda_*=8kA^{2N}+A.
\]
\end{lemma}

\begin{proof}
Let $\Omega$ be the diagonal matrix equal to $\rho^{(i-1)/2}$ on the coordinates
of gate $i$, and set $y=\Omega u$. Then $u^{\mathsf T}H^*u=y^{\mathsf T}\widetilde Hy$,
where $\widetilde H$ has diagonal blocks $H_i$ and off-diagonal entries coming
from the terms $-\rho^{i-1}u_{i,1}a_{ije}u_{j,e}=-\rho^{(i-j)/2}a_{ije}\,y_{i,1}y_{j,e}$
with $j<i$. The corresponding symmetric entries have modulus at most
$\frac12\rho^{1/2}\abs{a_{ije}}$. In the row of a coordinate $y_{i,1}$ these
entries come from its own radicand and from later gates that use it, and in the
row of $y_{j,e}$, $e\ge2$, only from later gates. Each family has absolute
coefficient sum at most $A$, so every absolute row sum of the off-diagonal part
is at most $\rho^{1/2}A=h_0/4$. Hence $\widetilde H\succeq\frac34h_0I$ by
Lemma~\ref{lem:reductions-local-exposing}(iii), and
$H^*=\Omega\widetilde H\Omega\succeq\frac34h_0\rho^{k-1}I$. Without rescaling, the
diagonal blocks contribute at most $\sum_i\rho^{i-1}\norm{H_i}\le8kA^{2N}$ and the
off-diagonal part, with entries at most $\frac12\abs{a_{ije}}$, at most $A$.
\end{proof}

The bound holds in every direction, for arbitrary signs of the coefficients,
for gates used by many later gates, and for radicands using several powers of
one predecessor.

\paragraph{A rational vanishing representation.}
For a monomial $\mathfrak m$ of degree at most two in the gate's own variables
$X_1,\ldots,X_n$, including $\mathfrak m=1$, let $w(\mathfrak m)$ be its weighted
degree, with $w(X_a)=a$, $w(X_aX_b)=a+b$ and $w(1)=0$, so that
$\mathfrak m(t,\ldots,t^n)=t^{w(\mathfrak m)}$. For $0\le h\le2n$ put $m_0=1$,
$m_h=X_h$ for $1\le h\le n$, and $m_h=X_nX_{h-n}$ for $n<h\le2n$. Let
$R_{\mathfrak m}=\mathfrak m-m_{w(\mathfrak m)}$. These rational quadratics
vanish on the power curve, so the gate-$i$ versions $R_{i,\mathfrak m}$ vanish at
$p$. Expanding~\eqref{eq:reductions-T-alpha} writes
$P_\alpha=\sum_{\mathfrak m}\chi_{\mathfrak m}(\alpha)\,\mathfrak m$ with
polynomials $\chi_{\mathfrak m}$ having integer coefficients, degrees at most
$2n$, and total absolute coefficient sum at most $8n$: each diagonal term
$\alpha^{d-1-2j}\ell_j^2$ and each neighboring pair
$\alpha^{d-2-2j}\ell_j\ell_{j+1}$ contributes at most four.

\begin{lemma}[Exact vanishing representation]
\label{lem:reductions-vanishing-basis}
With $S_i=X_{i,n_i}^2-b_i(X)X_{i,1}$ and $T_i=X_{i,n_i-1}X_{i,n_i}-b_i(X)=r_{i,n_i}$,
\[
 E_i^*=S_i-\alpha_iT_i+\sum_{\mathfrak m}\chi_{\mathfrak m}(\alpha_i)R_{i,\mathfrak m}.
\]
All of $S_i$, $T_i$ and $R_{i,\mathfrak m}$ are rational quadratics vanishing
at $p$.
\end{lemma}

\begin{proof}
Write $\alpha=\alpha_i$, $d=d_i$, $n=n_i$. Since
$\sum_{w(\mathfrak m)=h}\chi_{\mathfrak m}(\alpha)$ is the coefficient of $t^h$ in
$P_\alpha(t,\ldots,t^n)=t^{d+1}-\alpha t^d-\alpha^dt+\alpha^{d+1}$ and $d+1=2n$,
\[
 P_\alpha=\sum_{\mathfrak m}\chi_{\mathfrak m}(\alpha)R_{\mathfrak m}
 +X_n^2-\alpha X_{n-1}X_n-\alpha^dX_1+\alpha^{d+1}.
\]
Adding $(\alpha-X_1)(b_i-\alpha^d)=\alpha b_i-\alpha^{d+1}-X_1b_i+\alpha^dX_1$ gives
the identity. At $p$, $T_i=\alpha^d-\alpha^d=0$ and
$S_i=\alpha^{2n}-\alpha^d\alpha=0$.
\end{proof}

\paragraph{Assembly.}
Put
\[
 \mu=\gamma,\qquad\Lambda=\Lambda_*+1,\qquad
 \varepsilon_{\max}=\min\Bigl\{1,\frac{\mu^2}{2N},\frac{\nu^2\mu^2}{36N(\Lambda+N\beta)^2}\Bigr\},
\]
with $\beta,\nu$ from Lemma~\ref{lem:reductions-residuals}, and let $t=2^{-s}$
with $s$ the least nonnegative integer such that $\varepsilon=t^2\le\varepsilon_{\max}$. Let
\[
 B_0=A+1+32N^2(2A)^{2N},\qquad
 \tau=\min\Bigl\{1,\frac{\gamma}{2kB_0},\frac{\varepsilon}{2A^NkB_0}\Bigr\},
\]
compute rationals $\widehat\alpha_i$ with $\abs{\widehat\alpha_i-\alpha_i}\le\tau$
by Lemma~\ref{lem:reductions-root-approx}, and define
\[
 G=\sum_{i=1}^k\rho^{i-1}\Bigl[S_i-\widehat\alpha_iT_i
 +\sum_{\mathfrak m}\chi_{\mathfrak m}(\widehat\alpha_i)R_{i,\mathfrak m}\Bigr],
 \qquad
 F=\Bigl(\frac{G}{t\nu}\Bigr)^2+\sum_{i,e}\Bigl(\frac{r_{i,e}}{\nu}\Bigr)^2.
\]
By Lemma~\ref{lem:reductions-vanishing-basis}, $G$ is a rational quadratic with
$G(p)=0$, whatever the approximations. The difference $G-E^*$ has
$\norm{G-E^*}_1\le k\tau B_0$. Indeed, for each gate, $\abs{\widehat\alpha_i-\alpha_i}\norm{T_i}_1
\le\tau(A+1)$; and on $[-2A,2A]$ the derivative of $\chi_{\mathfrak m}$ is
at most its absolute coefficient sum times $2n(2A)^{2n-1}$, so
$\sum_{\mathfrak m}\abs{\chi_{\mathfrak m}(\widehat\alpha_i)-\chi_{\mathfrak m}(\alpha_i)}
\le16N^2(2A)^{2N}\tau$, while $\norm{R_{i,\mathfrak m}}_1\le2$. Since
$G-E^*$ vanishes at $p$ and $\norm p_\infty\le A^N$,
\eqref{eq:reductions-app-coefficients} gives
$G(p+u)=\ell^{\mathsf T}u+u^{\mathsf T}Hu$ with
\[
 \norm{H-H^*}\le k\tau B_0\le\tfrac12\gamma,\qquad
 \norm\ell\le2A^Nk\tau B_0\le\varepsilon .
\]
With Lemma~\ref{lem:reductions-coupled}, $\mu I\preceq H\preceq\Lambda I$.

Thus the hypotheses of Lemma~\ref{lem:quartic-realization} hold with $n=N$
variables, the $m=N$ residuals $r_{i,e}$, the quadratic $g=G$, the constants
$\mu,\Lambda,\beta,\nu$, and $\varepsilon=t^2$: the residual quadratic parts have
norm at most one and their gradients at $p$ have norm at most $\beta$
(Lemma~\ref{lem:reductions-residuals}); $\sum_{i,e}\nabla r_{i,e}(p)\nabla
r_{i,e}(p)^{\mathsf T}=J^{\mathsf T}J\succeq\nu^2I$; $G(p)=0$ with quadratic part
between $\mu I$ and $\Lambda I$ and $\norm{\nabla G(p)}\le\varepsilon$; and
$\varepsilon\le\varepsilon_{\max}$. The lemma shows that $F$ is a sum of $N+1$
squares of rational quadratics with $\nabla^2F\succeq\frac32I$ and unique zero
$p$, and that the canonical Hessian Gram of this representation is a rational
positive definite Hessian Gram of $F$, computed by polynomially many rational
operations.

\paragraph{Complexity.}
The numbers $A$, $h_0$, $\rho$, $\gamma=\mu$, $\Lambda$, $\beta$, $\nu$, $B_0$, $\varepsilon_{\max}$ and
$\tau$ are explicit rationals. Their bit lengths are polynomial in the
input length, because they are built from $A$ and $N\le L$ by products and powers
with exponents polynomial in $k$ and $N$; for example $\log_2(1/\nu)=(N-1)\log_2\beta$ and
$\log_2(1/\gamma)=O(kN\log_2A+k\log_2N)$. Lemma~\ref{lem:reductions-root-approx}
with $P=\lceil\log_2(1/\tau)\rceil$ runs in polynomial time. The polynomials
$\chi_{\mathfrak m}$ are explicit, there are $O(N^2)$ monomials per gate, and $F$
has $O(N^4)$ monomials. Hence the whole construction runs in polynomial time. The
normalization, the residuals, and the rational vanishing representation never
use a minimal polynomial of any gate value or of the joint field. \qed

\subsection{Proof of Theorem~\ref{thm:reductions-quaternion}}
\label{app:reductions-quaternion-compiler}

We use the notation of Section~\ref{sec:reductions-quaternion}, write
$\mathbf j=(0,0,1,0)$ and $\mathbf k=(0,0,0,1)$, and let
$e_1,e_2,e_3$ be the standard basis of the vector part. For unit quaternions put
$[u,w]=uw\bar u\bar w$, $\iota(u)=\mathbf iu\bar{\mathbf i}$,
$\operatorname{rot}(u)=\mathbf cu\bar{\mathbf c}$ with $\mathbf c=\frac12(1,1,1,1)$, $\operatorname{pr}(u)=u\,\iota(u)$,
$\operatorname{add}(u,w)=\operatorname{pr}(uw)$ and $\operatorname{prod}(u,w)=\operatorname{pr}(\operatorname{rot}([u,\operatorname{rot}(w)]))$.

\begin{lemma}[Quaternion identities]
\label{lem:reductions-quaternion-identities}
For unit quaternions $u,w$:
\begin{enumerate}[label=(\roman*)]
\item The three maps are
\[
 \iota(u)=(u_0,u_1,-u_2,-u_3),\qquad
 \operatorname{rot}(u)=(u_0,u_3,u_1,u_2),
\]
\[
 \operatorname{pr}(u)=(1-2u_1^2,\,2u_0u_1,\,2u_1u_3,\,-2u_1u_2).
\]
\item $[u,w]-\mathbf1=2(0,\vec u\times\vec w)\,\bar u\bar w$, hence
$\abs{[u,w]-\mathbf1}\le2\abs{\vec u}\abs{\vec w}$.
\item If $u_0,w_0\ge0$, then
\[
 \abs{[u,w]-\mathbf1-2(0,\vec u\times\vec w)}
 \le4\abs{\vec u}\abs{\vec w}(\abs{\vec u}+\abs{\vec w}).
\]
\item If $u_0\ge0$ and $\abs{\vec u}\le\frac12$, then
$\abs{\overrightarrow{\operatorname{pr}(u)}-2u_1e_1}\le4\abs{\vec u}^2$ and $\operatorname{pr}(u)_0\ge\frac12$.
\end{enumerate}
\end{lemma}

\begin{proof}
(i) Conjugation by a unit quaternion fixes the scalar part and acts linearly on
the vector part. Direct multiplication gives
$\mathbf i\mathbf j\bar{\mathbf i}=-\mathbf j$ and
$\mathbf i\mathbf k\bar{\mathbf i}=-\mathbf k$, which is the formula for $\iota$. For
$\operatorname{rot}$, write $\mathbf c=\frac12(\mathbf1+\mathbf i+\mathbf j+\mathbf k)$; expanding the
sixteen products in $(\mathbf c\mathbf i)\bar{\mathbf c}=\frac14(-\mathbf1+\mathbf i+
\mathbf j-\mathbf k)(\mathbf1-\mathbf i-\mathbf j-\mathbf k)$ gives $\mathbf j$. The
cyclic permutation $\mathbf i\mapsto\mathbf j\mapsto\mathbf k\mapsto\mathbf i$ is an
automorphism of the quaternions fixing $\mathbf c$, so also
$\mathbf c\mathbf j\bar{\mathbf c}=\mathbf k$ and $\mathbf c\mathbf k\bar{\mathbf c}=\mathbf i$.
Thus $\operatorname{rot}$ sends $u_1e_1+u_2e_2+u_3e_3$ to $u_1e_2+u_2e_3+u_3e_1$. Finally, for
$\vec u'=(u_1,-u_2,-u_3)$, the product $u\,\iota(u)$ has scalar part
$u_0^2-\vec u\cdot\vec u'=u_0^2-u_1^2+u_2^2+u_3^2=1-2u_1^2$, using $\abs u=1$, and
vector part $u_0(\vec u+\vec u')+\vec u\times\vec u'=(2u_0u_1,0,0)+(0,2u_1u_3,-2u_1u_2)$.

(ii) Since $wu\bar u\bar w=\mathbf1$, we have
$[u,w]-\mathbf1=(uw-wu)\bar u\bar w$, and $uw-wu=(0,2\vec u\times\vec w)$ by the
product formula. The norm is multiplicative and $\abs{\bar u}=\abs{\bar w}=1$.

(iii) The left side equals $\abs{2(0,\vec u\times\vec w)}\,\abs{\bar u\bar w-\mathbf1}$,
and $\abs{\bar u\bar w-\mathbf1}=\abs{u(\bar u\bar w-\mathbf1)}=\abs{\bar w-u}
\le\abs{w-\mathbf1}+\abs{u-\mathbf1}$. For a unit quaternion with $u_0\ge0$,
$u_0=(1-\abs{\vec u}^2)^{1/2}\ge1-\abs{\vec u}^2$, so
$\abs{u-\mathbf1}^2=2-2u_0\le2\abs{\vec u}^2$ and $\abs{u-\mathbf1}\le2\abs{\vec u}$.

(iv) By (i), $\overrightarrow{\operatorname{pr}(u)}-2u_1e_1=2u_1(u_0-1,u_3,-u_2)$. Here
$0\le1-u_0\le\abs{\vec u}^2\le\frac12\abs{\vec u}$, so the vector
$(u_0-1,u_3,-u_2)$ has norm at most $\frac{\sqrt5}2\abs{\vec u}\le2\abs{\vec u}$.
Also $\operatorname{pr}(u)_0=1-2u_1^2\ge\frac12$.
\end{proof}

\begin{lemma}[Small rational signal]
\label{lem:reductions-generator}
Let $g$ be a unit quaternion with $g_0>0$, $0<g_1\le2^{-16}$ and
$\abs{(g_2,g_3)}\le64g_1^2$. Then $g'=\operatorname{prod}(g,g)$ satisfies the same three
conditions and $g_1^2\le g_1'\le6g_1^2$. Consequently the iterates
$g^{(0)}=\mathbf w_\vartheta=\bigl(\frac{1-\vartheta^2}{1+\vartheta^2},
\frac{2\vartheta}{1+\vartheta^2},0,0\bigr)$, $\vartheta=2^{-20}$, and
$g^{(j+1)}=\operatorname{prod}(g^{(j)},g^{(j)})$ are rational unit quaternions with
$0<g^{(j)}_1<2^{-16\cdot2^j}$ and $\abs{(g^{(j)}_2,g^{(j)}_3)}\le64(g^{(j)}_1)^2$.
\end{lemma}

\begin{proof}
Write $x=g_1$. Then $\abs{\vec g}\le x+64x^2\le2x$. Let $U=\operatorname{rot}([g,\operatorname{rot}(g)])$. Both
$g$ and $\operatorname{rot}(g)$ have positive scalar parts and vector parts of norm
$\abs{\vec g}$. By Lemma~\ref{lem:reductions-quaternion-identities}(iii), and
since $\operatorname{rot}$ is an isometry fixing $\mathbf1$,
$\abs{U-\mathbf1-2(0,\operatorname{rot}(\vec g\times\operatorname{rot}(\vec g)))}\le8\abs{\vec g}^3\le64x^3$, where
$\operatorname{rot}$ acts on vectors by $(a_1,a_2,a_3)\mapsto(a_3,a_1,a_2)$. The first coordinate
of $\operatorname{rot}(\vec g\times\operatorname{rot}(\vec g))$ is $g_1^2-g_2g_3$, and
$2\abs{g_2g_3}\le g_2^2+g_3^2\le4096x^4\le x^3$. Hence
$\abs{U_1-2x^2}\le65x^3\le x^2$, so $x^2\le U_1\le3x^2$. By
Lemma~\ref{lem:reductions-quaternion-identities}(ii),
$\abs{U-\mathbf1}\le2\abs{\vec g}^2\le8x^2$, so $\abs{\vec U}\le8x^2$ and
$U_0\ge1-8x^2>\frac12$. By part~(i), $g'=\operatorname{pr}(U)$ has
$g'_1=2U_0U_1\in[x^2,6x^2]$, scalar part $1-2U_1^2>0$, and
$\abs{(g'_2,g'_3)}=2U_1\abs{(U_2,U_3)}\le48x^4\le64(g_1')^2$. Finally
$g'_1\le6x^2\le x\le2^{-16}$.

The constant $\mathbf w_\vartheta$ is a rational unit quaternion with positive
scalar part, $0<(\mathbf w_\vartheta)_1<2\vartheta=2^{-19}$ and zero transverse
part. Products and conjugates of rational unit quaternions are rational unit
quaternions. Iterating, $6g^{(j+1)}_1\le(6g^{(j)}_1)^2$, so
$6g^{(j)}_1\le(6g^{(0)}_1)^{2^j}<(2^{-16})^{2^j}$, and $g^{(j+1)}_1\ge(g^{(j)}_1)^2>0$.
\end{proof}

\paragraph{Quaternion signals.}
Fix $0<\delta<1$. A \emph{$(d,C,\Theta)$-signal} is now a unit quaternion $u$
with $u_0>0$, $\abs C\le\Theta$ and $\abs{\vec u-C\delta^de_1}\le\Theta\delta^{d+1}$,
where $d\ge1$, $C\in\Z$ and $\Theta\ge1$. Then $\abs{\vec u}\le2\Theta\delta^d$.

\begin{lemma}[Quaternion signal arithmetic]
\label{lem:reductions-quaternion-signals}
Suppose $\Theta\delta\le2^{-30}$ and let $\Theta'=2^{20}\Theta^4$.
\begin{enumerate}[label=(\roman*)]
\item If $u$ is a $(d,C,\Theta)$-signal, then $\bar u$ is a $(d,-C,\Theta)$-signal
and $\operatorname{pr}(u)$ is a $(d,2C,\Theta')$-signal.
\item If $u,w$ are $(d,C,\Theta)$- and $(d,C',\Theta)$-signals, then $\operatorname{add}(u,w)$ is a
$(d,2(C+C'),\Theta')$-signal.
\item If $u,w$ are $(a,C,\Theta)$- and $(b,C',\Theta)$-signals, then $\operatorname{prod}(u,w)$ is an
$(a+b,4CC',\Theta')$-signal.
\end{enumerate}
\end{lemma}

\begin{proof}
Let $h=2\Theta\delta^d$ in~(i) and~(ii); then $h\le2^{-29}$. All scalar parts
below are positive, and all vectors to which
Lemma~\ref{lem:reductions-quaternion-identities}(iv) is applied have norm at
most $\frac12$.

(i) Conjugation negates the vector part. For $\operatorname{pr}(u)$, part~(iv) of
Lemma~\ref{lem:reductions-quaternion-identities} gives
\[
 \abs{\overrightarrow{\operatorname{pr}(u)}-2C\delta^de_1}
 \le4h^2+2\Theta\delta^{d+1}\le18\Theta^2\delta^{d+1}.
\]

(ii) Write $\vec u=C\delta^de_1+\epsilon_u$ and $\vec w=C'\delta^de_1+\epsilon_w$.
Then
\[
 \overrightarrow{uw}-(C+C')\delta^de_1=\epsilon_w+\epsilon_u+(u_0-1)\vec w
 +(w_0-1)\vec u+\vec u\times\vec w.
\]
The first two terms have norm at most $2\Theta\delta^{d+1}$; $0\le1-u_0\le h^2$,
so the next two have norm at most $2h^3=16\Theta^3\delta^{3d}$; and since
$e_1\times e_1=0$, $\abs{\vec u\times\vec w}\le3\Theta^2\delta^{2d+1}$. As $d\ge1$, the
total is at most $21\Theta^3\delta^{d+1}$. Also $(uw)_0\ge u_0w_0-h^2>0$ and
$\abs{\overrightarrow{uw}}\le2h+h^2\le5\Theta\delta^d$. Applying
Lemma~\ref{lem:reductions-quaternion-identities}(iv) to $uw$,
$\abs{\overrightarrow{\operatorname{pr}(uw)}-2(C+C')\delta^de_1}\le4\cdot25\Theta^2\delta^{2d}+
42\Theta^3\delta^{d+1}\le142\Theta^3\delta^{d+1}$, and $\abs{2(C+C')}\le4\Theta$.

(iii) Now $\abs{\vec u}\le2\Theta\delta^a$ and $\abs{\vec w}\le2\Theta\delta^b$, and
$\operatorname{rot}(w)$ has positive scalar part and vector part $\operatorname{rot}(\vec w)=C'\delta^be_2+\operatorname{rot}(\epsilon_w)$.
Since $e_1\times e_2=e_3$,
$\abs{\vec u\times\operatorname{rot}(\vec w)-CC'\delta^{a+b}e_3}\le3\Theta^2\delta^{a+b+1}$. By
Lemma~\ref{lem:reductions-quaternion-identities}(iii),
$\abs{[u,\operatorname{rot}(w)]-\mathbf1-2(0,\vec u\times\operatorname{rot}(\vec w))}\le4\cdot4\Theta^2\delta^{a+b}
\cdot4\Theta\delta=64\Theta^3\delta^{a+b+1}$. The rotation $\operatorname{rot}$ maps $e_3$ to $e_1$,
so $U=\operatorname{rot}([u,\operatorname{rot}(w)])$ satisfies
$\abs{\vec U-2CC'\delta^{a+b}e_1}\le(6\Theta^2+64\Theta^3)\delta^{a+b+1}$, and by
part~(ii) of that lemma $\abs{U-\mathbf1}\le8\Theta^2\delta^{a+b}$, so
$U_0>0$ and $\abs{\vec U}\le8\Theta^2\delta^{a+b}$. Applying part~(iv),
$\abs{\overrightarrow{\operatorname{pr}(U)}-4CC'\delta^{a+b}e_1}\le256\Theta^4\delta^{2(a+b)}
+(12\Theta^2+128\Theta^3)\delta^{a+b+1}\le396\Theta^4\delta^{a+b+1}$, and
$\abs{4CC'}\le4\Theta^2$. All bounds are at most $\Theta'$ times the relevant
power of $\delta$.
\end{proof}

As in the cube-root compiler, no bound divides by a coefficient, so zero
values, cancellation, and repeated arguments such as $\operatorname{prod}(u,u)$ or $\operatorname{add}(u,\bar u)$
are covered.

\paragraph{The compiler.}
Prepare the program as in Appendix~\ref{app:reductions-root-compiler}, so that
its output is $W=2V-1$. Each node $\nu$ receives a pair of quaternion signals
$(f_\nu,g_\nu)$ of a common order $d_\nu$ and prescribed coefficients
$(C_\nu,D_\nu)$ with $D_\nu>0$ and $C_\nu=D_\nu V_\nu$. The constant $1$ receives
$(g,g)$ and the constant $0$ receives $(\mathbf1,g)$, both of order one, where
$g$ is the generator quaternion fixed below, with coefficients $(1,1)$ and
$(0,1)$. A product $\nu=\nu_1\nu_2$ receives $f_\nu=\operatorname{prod}(f_{\nu_1},f_{\nu_2})$ and
$g_\nu=\operatorname{prod}(g_{\nu_1},g_{\nu_2})$, with coefficients $(4C_{\nu_1}C_{\nu_2},
4D_{\nu_1}D_{\nu_2})$. A sum or difference receives
\[
 f_\nu=\operatorname{add}\bigl(\operatorname{prod}(f_{\nu_1},g_{\nu_2}),\,\operatorname{prod}(f_{\nu_2},g_{\nu_1})^{\pm1}\bigr),\qquad
 g_\nu=\operatorname{pr}\bigl(\operatorname{prod}(g_{\nu_1},g_{\nu_2})\bigr),
\]
where $u^{+1}=u$ and $u^{-1}=\bar u$, with coefficients
$(8(C_{\nu_1}D_{\nu_2}\pm C_{\nu_2}D_{\nu_1}),\,8D_{\nu_1}D_{\nu_2})$. In both cases
$d_\nu=d_{\nu_1}+d_{\nu_2}$, the two arguments of $\operatorname{add}$ have equal order, and
$C_\nu=D_\nu V_\nu$ with $D_\nu>0$. Let $T$ be the number of signal operations
$\operatorname{prod}$, $\operatorname{add}$, $\operatorname{pr}$ and conjugation
used.

Let $r_0=2T+2$ and let $g=g^{(r_0)}$ from Lemma~\ref{lem:reductions-generator},
with $\delta=g_1$. Then $g$ is a $(1,1,128)$-signal and $\mathbf1$ a
$(1,0,128)$-signal. Put $\Theta_0=128$ and $\Theta_t=2^{20}\Theta_{t-1}^4$, so that
$\log_2\Theta_t=(41\cdot4^t-20)/3<14\cdot4^t$. Since
$\delta<2^{-16\cdot2^{2T+2}}=2^{-64\cdot4^T}$, we get
$\Theta_T\delta<2^{-50\cdot4^T}\le2^{-30}$. Induction over the operations, using
Lemma~\ref{lem:reductions-quaternion-signals}, shows that every signal is a
$(d,C,\Theta_T)$-signal for its order and prescribed coefficient, for any
sharing pattern. For the output node, $C=D\,W$ is a nonzero integer, so the
selected coordinate $x$ of $f_\nu$ satisfies
$\abs{x-C\delta^d}\le\Theta_T\delta^{d+1}<\frac12\delta^d\le\frac12\abs C\delta^d$.
Hence $x\ne0$ and $\operatorname{sign}x=\operatorname{sign}W$.

Each of $\iota$, $\operatorname{rot}$, $\operatorname{pr}$, $\operatorname{add}$, $\operatorname{prod}$ and the commutator is a fixed composition of at
most twelve products and conjugations with the constants $\mathbf i,\mathbf c$
and their conjugates, so the circuit has $O(T)$ gates and is computed in
polynomial time. Every gate value is a rational unit quaternion, and nothing
is expanded.

For the upper bound, first check by rational arithmetic that every given
constant has norm one, and map malformed inputs to a fixed program with output
$-1$. Then represent each gate value $u$ by four integer straight-line programs
$N_0,\ldots,N_3$ and one program $D>0$ with $u_s=N_s/D$; a constant with
binary numerators and denominators has a representation of linear size. A product multiplies
the denominators and forms the four bilinear numerators of the product formula;
a conjugate negates three numerators. The sign of a coordinate $u_s$ is the
sign of $N_s$, so $u_s>0$ is one $\PosSLP$ instance, and $u_s\ge0$ is the
instance $2N_s+1>0$. \qed

\subsection{Proof of Proposition~\ref{prop:reductions-quaternion-realization}}
\label{app:reductions-quaternion-realization}

Let the circuit have gates $1,\ldots,k$ with values $p_i$, and variables
$X_i\in\R^4$. For any quaternions, $\abs{ab}=\abs a\abs b$,
$\langle ab,c\rangle=\langle a,c\bar b\rangle$, and left or right multiplication
by a unit quaternion is an orthogonal map of $\R^4$.

\paragraph{Residuals.}
The vector residual of gate $i$ is $r_i=X_i-p_i$ for a constant gate, whose value $p_i$ is given,
$r_i=X_i-X_aX_b$ for a product, and $r_i=X_i-\bar X_a$ for a conjugate. Its four
components are rational quadratics. Their common real zero is exactly $p$, by
induction over the gates. At $p$ the Jacobian $J$ is block lower triangular
with identity diagonal blocks, so $\det J=1$. The block row of a product gate
is $(I,-\mathcal R_{p_b},-\mathcal L_{p_a})$, with $\mathcal R_{p_b}\colon x\mapsto xp_b$ and
$\mathcal L_{p_a}\colon x\mapsto p_ax$ orthogonal; if $a=b$ the parent block is
$-(\mathcal R_{p_a}+\mathcal L_{p_a})$. The block row of a conjugate gate is
$(I,-\diag(1,-1,-1,-1))$. Hence every block row has operator norm at most three,
every scalar residual gradient has norm at most three, and
$\norm J_F^2\le\sum_i4\cdot9=9N$. With $\beta=4N$ and $\nu=\beta^{-(N-1)}$,
\eqref{eq:reductions-app-singular} gives $J^{\mathsf T}J\succeq\nu^2I$. A
component of a product $X_aX_b$ with $a\ne b$ is a bilinear form
$X_a^{\mathsf T}KX_b$ with $K$ a signed permutation matrix, whose symmetric
matrix has norm $\frac12$. If $a=b$ and $X_a=(x_0,x_1,x_2,x_3)$, the components are
$x_0^2-x_1^2-x_2^2-x_3^2$, $2x_0x_1$, $2x_0x_2$ and $2x_0x_3$, whose matrices have norm one. Constant and
conjugate residuals are affine. So every residual quadratic part has norm at
most one.

\paragraph{Exposing form.}
Let $q_i=\abs{X_i}^2-1$ and define
\[
 E_i=q_i-2\langle p_i,r_i\rangle-
 \begin{cases}0&\text{constant gate},\\ q_a+q_b&\text{product gate},\\
 q_a&\text{conjugate gate},\end{cases}
\]
where for $a=b$ the product term is $2q_a$. For a product gate, expanding
and using $\langle X_aX_b,p_i\rangle=\langle X_a,p_i\bar X_b\rangle$ and
$\abs{p_i\bar X_b}=\abs{X_b}$ gives
\[
 E_i=\abs{X_i-p_i}^2-\abs{X_a-p_i\bar X_b}^2 .
\]
Since $p_i=p_ap_b$ and $\abs{p_b}=1$, we have $p_i\bar p_b=p_a$, so with
$U=X-p$,
\begin{equation}
 E_i(p+U)=\abs{U_i}^2-\abs{U_a-p_i\bar U_b}^2 .
 \label{eq:reductions-quaternion-exposing}
\end{equation}
The order of the factors matters because multiplication is not commutative,
and~\eqref{eq:reductions-quaternion-exposing} holds verbatim when $a=b$. For a
conjugate gate, $p_i=\bar p_a$ gives $E_i(p+U)=\abs{U_i}^2-\abs{U_a}^2$; for a
constant gate, $E_i(p+U)=\abs{U_i}^2$. In each case $E_i$ has zero value and
zero gradient at $p$, and its negative part is at most
$2\abs{U_a}^2+2\abs{U_b}^2$, or $4\abs{U_a}^2$ if $a=b$, or $\abs{U_a}^2$.

Let $\omega_i=16^{-(i-1)}$, $\omega_*=\omega_k$ and $E^*=\sum_i\omega_iE_i$, and write
$E^*(p+U)=U^{\mathsf T}H_*U$. Every block $U_j$ receives $\omega_j\abs{U_j}^2$ from
its own gate and loses at most $4\omega_i\abs{U_j}^2$ to each later gate $i$,
whatever the sharing pattern. As $\sum_{i>j}\omega_i\le\omega_j/15$,
\[
 \tfrac{11}{15}\,\omega_*I\preceq H_*\preceq I .
\]

\paragraph{Rounding.}
Fix $\varepsilon=t^2$ below and let
$\eta=\min\{1,\omega_*/(32k),\varepsilon/(16k)\}$. Choose a dyadic mesh
$h\le\eta/(2\cdot3^k)$ and compute $\widehat p_i$ in topological order:
$\widehat p_i$ rounds each coordinate of the given constant $p_i$, of $\widehat p_a\widehat p_b$, or of
$\overline{\widehat p_a}$ to the nearest multiple of $h$. Let
$e_i=\abs{\widehat p_i-p_i}$. Rounding changes a vector by at most $h$, and if
$e_a,e_b\le1$ then
$\abs{\widehat p_a\widehat p_b-p_ap_b}\le e_a(1+e_b)+e_b\le3\max\{e_a,e_b\}$.
Induction gives $e_i\le h(3^i-1)\le\eta/2$; in particular every coordinate error
is at most $\eta$, and the assumption $e\le1$ is maintained. The mesh has
$O(k+\log_2(1/\eta))$ bits, and the gate fractions are never expanded. Define
$G$ as $E^*$ with each coefficient vector $p_i$ in the term
$-2\langle p_i,r_i\rangle$ replaced by $\widehat p_i$, and with every residual and
every $q_i$ kept exact. Then $G(p)=0$. Writing $G(p+U)=\ell^{\mathsf T}U+U^{\mathsf T}HU$
and $\delta_i=\widehat p_i-p_i$, the difference
$G-E^*=-2\sum_i\omega_i\langle\delta_i,r_i\rangle$ gives
$\norm\ell\le2\sum_i3\abs{\delta_i}\le12k\eta$ and
$\norm{H-H_*}\le2\sum_i\norm{\delta_i}_1\le8k\eta$. With
$\mu=\omega_*/4$ and $\Lambda=2$, the choice of $\eta$ gives
\[
 \mu I\preceq H\preceq\Lambda I,\qquad\norm{\nabla G(p)}=\norm\ell\le\varepsilon .
\]

\paragraph{Assembly.}
Let $t=2^{-s}$ with $s$ the least nonnegative integer such that
\[
 \varepsilon=t^2\le\min\Bigl\{1,\frac{\mu^2}{2N},\frac{\nu^2\mu^2}{36N(\Lambda+N\beta)^2}\Bigr\},
\]
and let
\[
 F=\Bigl(\frac{G}{t\nu}\Bigr)^2+\frac1{\nu^2}\sum_{i=1}^k\abs{r_i}^2,
\]
where $\abs{r_i}^2$ is the sum of the squares of the four components of $r_i$.
The hypotheses of Lemma~\ref{lem:quartic-realization} hold with $n=N$
variables, the $m=N$ scalar residuals, $g=G$, and the constants
$\mu,\Lambda,\beta,\nu,\varepsilon$. Hence $F$ is a sum of $N+1$ squares of
rational quadratics with $\nabla^2F\succeq\frac32I$ and unique zero $p$, and the
canonical Hessian Gram of this representation is a rational positive definite
Hessian Gram of $F$, computed in polynomial time. The derived scale and
rounding parameters $\omega_i$, $\mu$, $\Lambda$, $\beta$, $\nu$,
$\varepsilon$, $t$, $\eta$, and $h$, and the rounded coefficients
$\widehat p_i$, have $O(k+N\log N)$ bits. The given rational generators may
have coordinates of any bit length; they enter the constant-gate residuals
exactly. Every coefficient of $G$, of the square factors of $F$, and of its
Hessian Gram therefore has bit length polynomial in the total input length.
The zero is the exact list of gate values, so $p\in\Q^N\cap[-1,1]^N$. \qed

\subsection{Proof of Proposition~\ref{prop:reductions-box-baseline}}
\label{app:reductions-comparisons}

Every $y_t$ has a triangular bound $y_t\le g_t(y_{<t})$: for a root variable,
$g_t=\rho_i(Y)^{1/d_i}$, where the radicand is affine in the earlier power
variables and is at least $\underline\xi_i^{\,d_i}\ge1$ on the boxes by the direct
interval test; for a power variable, $g_t=-X_i^e$. The feasible set is convex,
because powers are convex on the positive boxes, and compact. The circuit point
$y^*$ is feasible with every bound tight. On the boxes, which are convex, the
partial derivatives of $g_t$ are bounded by $\abs{a_{ije}}\le A$ for a root and by
$eX^{e-1}\le NA^{N-1}$ for a power, so
$\abs{g_t(y)-g_t(y')}\le C\sum_{s<t}\abs{y_s-y'_s}$ with $C=NA^N$.

Let $y$ be feasible, $\delta_t=g_t(y_{<t})-y_t\ge0$ and $\epsilon_t=y^*_t-y_t$. Then
$\epsilon_t=\delta_t+g_t(y^*_{<t})-g_t(y_{<t})$, so
$\abs{\epsilon_t}\le\delta_t+C\Sigma_{t-1}$ with $\Sigma_t=\sum_{s\le t}\abs{\epsilon_s}$.
Hence $\Sigma_t\le(1+C)\Sigma_{t-1}+\delta_t$ and
$\Sigma_{t-1}\le\sum_{s<t}(1+C)^{t-1-s}\delta_s$. With $B=2(C+1)$ and $W_t=B^{N-t}$,
\[
 \sum_tW_t\epsilon_t\ge\sum_tW_t(\delta_t-C\Sigma_{t-1})
 \ge\sum_s\delta_s\Bigl[W_s-C\sum_{t>s}W_t(1+C)^{t-1-s}\Bigr]
 \ge\frac1{1+C}\sum_sW_s\delta_s,
\]
because $\sum_{t>s}W_t(1+C)^{t-1-s}\le W_s\sum_{h\ge1}B^{-h}(1+C)^{h-1}
=W_s/(B-1-C)=W_s/(1+C)$. Thus $\sum_tW_ty^*_t\ge\sum_tW_ty_t$, with equality only
if every $\delta_s=0$, in which case triangular recursion gives $y=y^*$. Finally,
$\log_2W_t\le N\log_2(2C+2)$ is polynomial, and all other data are given. \qed
